\scanpage{108}

on a donc
\[
(8)\qquad \mathcal E_X^{\unsure{\xi|\zeta}}\longrightarrow\mathcal E_Y^\eta
\]
si on a un hom de schémas ``\unsure{géom\textsuperscript{t}.} \unsure{ponctués|pointés}'' $(X,\unsure{\xi|\zeta})\longrightarrow(Y,\eta)$. On notera que \unsure{tte|tt} ponctuation géom. de $X$ en un $x\in X$ --- i.e. une ext. sep. close $\Omega$ de $k(x)$ --- définit \inserted{\unsure{can\textsuperscript{t}}} une ponctuation \inserted{géom.} aussi de $Y$ en $y=f(x)$, \deleted{par} par composition --- $\Omega$ ``est'' une extension de $k(y)$ puisque $\unsure{k(x)}\hookleftarrow k(y)$ --- et l'hom (8) s'identifie aussi à
\[
(9)\qquad \mathcal E_X^{\overline{k(x)}}\longrightarrow\mathcal E_Y^{\overline{k(y)}},
\]
où $\overline{k(x)}$, $\overline{k(y)}$ sont les clôtures séparables \uline{dans $\Omega$}.

On posera
\[
(10)\qquad \mathcal E_{X/Y}^{\widetilde X}=\operatorname{Ker}\bigl(\mathcal E_X^{\widetilde X}\longrightarrow\mathcal E_Y^{\widetilde X}=\mathcal E_Y^{\widetilde Y}\bigr)
\]

\margin{107}
